\section{Phase I: Data Integrity Verification and Cleaning}
\label{sec:phase1_cleaning}

\subsection{Sanity checks}
\label{sec:phase1_sanity}

We verify:
\begin{itemize}[leftmargin=*]
  \item Window integrity: each project has seven consecutive windows from eight frozen tags.
  \item Arithmetic: for every row, \texttt{churn} $=$ \texttt{added} $+$ \texttt{deleted}.
  \item Totals: window totals equal sums of developer rows.
\end{itemize}

All 28 windows pass integrity checks, and no window has zero commits or zero churn in the cleaned totals.

\subsection{Canonical identity}
\label{sec:phase1_dedup}

The regenerated extraction pipeline canonicalizes developer identity during extraction. The canonical aggregation key is:
\[
(\texttt{project\_id}, \texttt{window\_id}, \texttt{author\_email})
\]
The extractor aggregates commits by email and selects a canonical \texttt{author\_name} using the most frequent non-\texttt{unknown} value.
This produces \texttt{metrics\_window\_dev.csv} with 273 developer-window rows. The cleaning pass keeps the same
273 canonical rows and records 0 duplicate key groups merged.

\subsection{Bot labeling}
\label{sec:phase1_bots}

We retain an \texttt{is\_bot} flag to allow reporting with and without automated accounts. In the Initial Cohort, 18 developer-window rows
are flagged as bots in the cleaned dataset. The pipeline now exports both all-account and human-only analysis populations:
all-account outputs remain the audit trail for total repository activity, while human-only outputs are the primary source for
M3 cadence and M4 ownership/concentration. This prevents automated dependency-update or formatting accounts from being interpreted
as human maintainers.

\subsection{Known limitation: commit-date/tag-date drift}
\label{sec:phase1_author_date_lim}

Some commits included between tags carry committer timestamps that precede the tag-date window boundaries.
This can happen when commits were authored or committed on side branches before they became reachable in the tagged history.
This affects optional cadence features based on timestamps (first/last commit), but does not affect churn or commit counts
computed from the Git revision range. Phase II therefore treats human-only commit counts and duration-normalized rates as primary cadence evidence,
and uses timestamp-derived active-day features only with sensitivity checks.

\subsection{Issue-resolution ledger}
\label{sec:phase1_issue_ledger}

\begin{table}[H]
\centering
\renewcommand{\arraystretch}{1.15}
\begin{tabularx}{\textwidth}{p{4.2cm}p{2.2cm}X}
\toprule
\textbf{Issue} & \textbf{Observed count} & \textbf{Initial Cohort resolution} \\
\midrule
Display-name variation for the same email/window & 0 groups after regeneration & Canonical aggregation by email key is performed during extraction. \\
Incorrect unique-author counting in raw totals & 0 windows after regeneration & Raw and clean totals compute unique authors by \texttt{author\_email}. \\
Bot contamination in ownership shares & 18 rows & Preserved \texttt{is\_bot}, exported human-only M3/M4 outputs, and generated bot-summary CSVs for sensitivity analysis. \\
Commit timestamp outside tag-date boundaries & 45 rows & Documented limitation; Phase II sensitivity checks avoid over-reliance on timestamp-derived cadence fields. \\
\bottomrule
\end{tabularx}
\caption{Initial Cohort issue ledger and implemented cleaning/validation actions.}
\label{tab:round1_issue_ledger}
\end{table}





